\chapter{Theory of  Production}

\section{Primal Representation of Technology}

Representation of technology can be expressed in several equivalent ways:

\begin{enumerate}
    \item Production possibilities set
    \item Input requirement set
    \item Producible-output set
    \item Production (transformation) function
    \item Distance function
\end{enumerate}


\subsection{Production Possibilities Set $T$}

\[
T = \{(x,y) : x \text{ can produce } y\},
\qquad x \in \mathbb{R}^N_+,\; y \in \mathbb{R}^M_+.
\]

\subsection*{Properties of $T$}

\begin{itemize}
    \item[T.1] $T$ is nonempty.
    
    \item[T.2] $T$ is closed.
    
    \item[T.3] $T$ is convex.
    
    \item[T.4] If $(x,y) \in T$ and $x' \ge x$, then $(x',y) \in T$
    \\(free disposability of inputs).
    
    \item[T.5] If $(x,y) \in T$ and $y' \le y$, then $(x,y') \in T$
    \\(free disposability of outputs).
    
    \item[T.6] For every finite $x$, the set of feasible outputs is bounded above.
    
    \item[T.7] $(x,0_M) \in T$, but if $y \ge 0$, then $(0_N,y) \notin T$
    \\(weak essentiality / no free lunch).
\end{itemize}

\medskip

\textbf{Interpretation of properties:}

\begin{itemize}
    \item[T.1] There exist feasible input-output combinations.
    
    \item[T.2] The boundary of $T$ is included in $T$ (no holes).
    
    \item[T.3] If $(x,y) \in T$ and $(x',y') \in T$, then
    \[
    (\alpha x + (1-\alpha)x',\ \alpha y + (1-\alpha)y') \in T,
    \quad \alpha \in [0,1].
    \]
    This is closely related to the concept of a diminishing marginal
    rate of technical substitution (MRTS).

    \item[T.4] If an input bundle $x$ can produce a given output bundle $y$,
    then any larger input bundle $x' \ge x$ can also produce $y$.

    \item[T.5] If $(x,y) \in T$, then producing a smaller output bundle
    $y' \le y$ does not require more input.

    \item[T.6] This is a technical condition ensuring the existence of
    a production function and guaranteeing that optimization problems
    have solutions.

    \item[T.7] This represents the \textit{no free lunch} condition:
    zero output is always feasible, but positive output requires
    positive input.
\end{itemize}

\subsection{Input Requirement Set $V(y)$}

\[
V(y)
=
\{x : (x,y) \in T\},
\]
that is, the set of all input bundles capable of producing
the output bundle $y$.

\subsection*{Properties of $V(y)$}

\begin{itemize}
    \item[V.1] $V(y)$ is nonempty for at least one finite $y$.

    \item[V.2] $V(y)$ is closed.

    \item[V.3] $V(y)$ is convex.

    \item[V.4] If $x \in V(y)$ and $x' \ge x$, then $x' \in V(y)$
    \\(free disposability of inputs).

    \item[V.5] If $y' \ge y \ge 0$, then
    \[
    V(y') \subseteq V(y).
    \]

    \item[V.6] $0_N \notin V(0)$ for $y \ge 0_M$, but $0_N \in y(0_M)$
    \\(follows from T.7: zero output is always producible).
\end{itemize}

\subsection*{Proof of properties of $V(y)$}

$V(y)$ is nonempty for at least one finite output vector since T.1 implies
that at least one feasible input-output combination exists.

However, it does not imply that $V(y)$ is nonempty for all $y$.

\medskip

V.2 follows directly from T.2.

\medskip

V.3: Suppose $x,x' \in V(y)$. Then $(x,y) \in T$ and $(x',y) \in T$.
By convexity of $T$,
\[
(\alpha x + (1-\alpha)x',\ \alpha y + (1-\alpha)y)
=
(\alpha x + (1-\alpha)x',\ y) \in T.
\]
Thus,
\[
\alpha x + (1-\alpha)x' \in V(y),
\]
so $V(y)$ is convex.

\medskip

V.4 and V.5 follow from T.4 and T.5, respectively.

\medskip

T.7 implies V.6. \hfill $\blacksquare$

\subsection*{Convexity and Production Functions}

The conditions V.3 and T.3 are closely related but not equivalent.

\medskip

V.3 implies that the set
\[
V(y)=\{x : f(x) \ge y\}
\]
is convex, and thus $f(x)$ is a \textbf{quasiconcave function}.

\medskip

\subsection*{Definition 1 (Concavity)}

Let $f$ be a function defined on a convex set $A \subset \mathbb{R}^N$.

Define
\[
T = \{(x,y) \in \mathbb{R}^{N+1} : x \in A,\ y \le f(x)\}.
\]

Then $f$ is \textbf{concave} on $A$ if and only if $T$ is convex.

\medskip

\subsection*{Implications}

\begin{itemize}
    \item T.3 $\Rightarrow$ $f(x)$ is concave.
    \item V.3 $\Rightarrow$ $f(x)$ is quasiconcave.
\end{itemize}

Thus, T.3 is more restrictive than V.3.



\subsection{Producible-Output Set $Y(x)$}

\[
Y(x)=\{y : (x,y)\in T\},
\]
that is, the set of all output bundles that can be produced using
a given input bundle $x$.

\subsection*{Properties of $Y(x)$}

\begin{itemize}
    \item[Y.1] $Y(x)$ is nonempty.
    
    \item[Y.2] $Y(x)$ is closed.
    
    \item[Y.3] $Y(x)$ is convex.
    
    \item[Y.4] If $x' \ge x$, then
    \[
    Y(x') \supseteq Y(x).
    \]
    
    \item[Y.5] If $y \in Y(x)$ and $y' \le y$, then
    \[
    y' \in Y(x).
    \]
    
    \item[Y.6] $Y(x)$ is bounded from above for finite $x$.
    
    \item[Y.7] If $y \ne 0$, then
    \[
    y \notin Y(0_N),
    \qquad
    0_M \in Y(x).
    \]
\end{itemize}

\begin{tcolorbox}[breakable]

\subsection*{Exercise 2-1}


Prove Y.1--Y.7.

\end{tcolorbox}

\medskip

The frontiers of $T$, $V(y)$, and $Y(x)$ are, respectively,


\begin{itemize}
    \item the production function,
    \item the isoquant,
    \item the production possibility frontier.
\end{itemize}

Thus, $T$, $V(y)$, and $Y(x)$ convey almost identical information
about production technology.




\subsection*{Example 1 (CRS Technology)}

Representing CRS technology:


The technology set $T$ exhibits \textbf{constant returns to scale (CRS)} if
\[
(x,y) \in T \;\Rightarrow\; (\theta x,\theta y) \in T,
\qquad \theta > 0.
\]

\medskip

Under CRS,
\[
\begin{aligned}
V(\theta y)
&=
\{x : (x,\theta y)\in T\}\\
&=
\{x : (x/\theta, y)\in T\}\\
&=
\theta \{x/\theta : (x/\theta,y)\in T\}\\
&=
\theta V(y).
\end{aligned}\]

Similarly,
\[
\begin{aligned}
Y(\theta x)
&=
\{y : (\theta x,y)\in T\}\\
&=
\{y : (x, y/\theta)\in T\}\\
&=
\theta \{y/\theta : (x,y/\theta)\in T\}\\
&=
\theta Y(x).
\end{aligned}\]

\subsection{Production or Transformation Function}

\subsubsection{A Production Function}

The production function is defined as
\[
\begin{aligned}
f(x)
&=
\max \{y : (x,y)\in T\}\\
&=
\max \{y \in Y(x)\},\\
&y \in \mathbb{R}_+.
\end{aligned}
\]

\subsection*{Properties of the Production Function}

\begin{itemize}
    \item[f.1] $f(x)$ exists for all finite $x$.
    
    \item[f.2] If $x' \ge x$, then
    \[
    f(x') \ge f(x)
    \quad \text{(monotonicity)}.
    \]
    
    \item[f.3] $f(x)$ is quasiconcave.
    
    \item[f.4] $f(0_N)=0$
    \quad \text{(weak essentiality)}.
\end{itemize}

\subsection*{Proof}

From T.2 and T.6, the set $T$ is compact.
By the Weierstrass theorem, the maximum exists, so $f(x)$ is well-defined.

\medskip

From T.4, if $x' \ge x$, then
\[
Y(x) \subseteq Y(x'),
\]
so the feasible set expands, implying
\[
\max Y(x') \ge \max Y(x),
\]
hence $f(x') \ge f(x)$.

\medskip

Since
\[
V(y) = \{x : f(x) \ge y\}
\]
is convex, it follows that $f(x)$ is quasiconcave.
\hfill $\blacksquare$






\subsubsection{A Transformation Function}

A transformation function is an implicit representation of technology.

For multi-output production, define
\[
Y(x,y)=0, \qquad y \in \mathbb{R}^M_+.
\]

For the scalar-output case,
\[
Y(x,y)=y - f(x)=0.
\]

\medskip

Note that
\[
\begin{aligned}
&V(y)=\{x : Y(x,y)\le 0\},\\
&Y(x)=\{y : Y(x,y)\le 0\}.
\end{aligned}
\]

\subsection*{Properties of the Transformation Function}

\begin{itemize}
    \item[J.1] If $Y(x_0,y_0)\le 0$ and $Y(x,y)\le 0$, then
    \[
    \begin{aligned}
    &Y(\theta x_0+(1-\theta)x,\ \theta y_0+(1-\theta)y)\\
    &\le
    \theta Y(x_0,y_0)+(1-\theta)Y(x,y)\\
    &(\text{convexity}).
    \end{aligned}
    \]

    \item[J.2] If $Y(x,y)\le 0$, then
    \[
    Y(x',y)\le 0 \quad \text{for } x' \ge x.
    \]

    \item[J.3] 
    \[
    Y(0_N,y)>0 \quad \text{for } y \ne 0_M.
    \]
\end{itemize}

\subsection{Output and Input Distance Functions}

\subsubsection{Output Distance Function}

\[
D_o(x,y)
=
\inf \{\theta > 0 : y/\theta \in Y(x)\}.
\]

\medskip

The output distance function measures the maximum proportional expansion
of the output vector $y$ that is feasible given input $x$.

That is, it pushes $y$ as far as possible away from the origin along
the ray through $y$, while remaining in the feasible set $Y(x)$.


\subsection*{Output Distance Function}

The value of the output distance function at point $a$ for $b$ on the edge is $\frac{0a}{0b}$.
Notice that
\[
y \in Y(x) \;\Longleftrightarrow\; D_o(x,y) \le 1.
\]


\begin{itemize}
    \item[$D_o$.1] $D_o(x,0)=0$, and $D_o(0,y)=+\infty \;\; \text{if } y \ne 0.$

    \item[$D_o$.2] $D_o(x',y) \le D_o(x,y)$ if $x \le x'$\\
    \text{(non-increasing in $x$)}.
    
    \item[$D_o$.3] $D_o(x,y) \ge D_o(x,y') \quad \text{if } y' \le y $\\
    \text{(non-decreasing in $y$)}.
    
    \item[$D_o$.4] $D_o(x,\mu y) = \mu D_o(x,y), \quad \mu > 0$\\
    \text{(homogeneous of degree 1 in $y$)}.
    
    \item[$D_o$.5] $D_o\bigl(x, \lambda y + (1-\lambda)y'\bigr) \le \lambda D_o(x,y) + (1-\lambda) D_o(x,y'),$
    $\quad 0 \le \lambda \le 1$
    $\quad \text{(convex in $y$)}.$
\end{itemize}

\subsubsection*{Proof}

First,
\[
D_o(x,0_M)
=
\inf \{\theta > 0 : 0_M \in Y(x)\}
= 0,
\]
since $0_M \in Y(x)$.

If $y \ne 0$, then $y \notin Y(0_N)$, hence
\[
D_o(0_N,y) = +\infty.
\]

By monotonicity of the production set (Y.4),
\[
x' \ge x \;\Rightarrow\; Y(x') \supseteq Y(x),
\]
so
\[
D_o(x',y) \le D_o(x,y).
\]

Monotonicity in outputs ($D_o$.3) follows from free disposability (Y.5).

For homogeneity,
\[
D_o(x,\mu y)
=
\inf \{\theta > 0 : \mu y/\theta \in Y(x)\}.
\]
Let $\theta = \mu t$. Then
\[
D_o(x,\mu y)
=
\mu \inf \{t > 0 : y/t \in Y(x)\}
=
\mu D_o(x,y).
\]

Before we prove $D_o$.5, we first show that $D_o(x,y)$ is subadditive in outputs:
\[
D_o(x,y+y') \le D_o(x,y) + D_o(x,y').
\]

From $D_o$.4, we know that
\[
D_o\left(x,\frac{y}{D_o(x,y)}\right)
=
\frac{1}{D_o(x,y)} D_o(x,y)
=1.
\]

Therefore,
\[
\frac{y}{D_o(x,y)} \in Y(x),
\qquad
\frac{y'}{D_o(x,y')} \in Y(x).
\]

Because $Y(x)$ is convex, it follows that
\[
\lambda \frac{y}{D_o(x,y)}
+
(1-\lambda)\frac{y'}{D_o(x,y')}
\in Y(x),
\qquad 0 < \lambda < 1.
\]

Let
\[
\lambda = \frac{D_o(x,y)}{D_o(x,y)+D_o(x,y')}.
\]

Then,
\[
\frac{y+y'}{D_o(x,y)+D_o(x,y')} \in Y(x).
\]

Therefore,
\[
D_o\left(x,\frac{y+y'}{D_o(x,y)+D_o(x,y')}\right) \le 1.
\]

By $D_o$.4, this implies
\[
D_o(x,y+y') \le D_o(x,y) + D_o(x,y').
\]

Finally, take $y=\lambda y$ and $y'=(1-\lambda)y'$. Then
\[
D_o(x,y+y')
=
D_o\bigl(x,\lambda y + (1-\lambda)y'\bigr)
\]
\[
\le
D_o(x,y) + D_o(x,y')
=
\lambda D_o(x,y) + (1-\lambda)D_o(x,y').
\]




\subsubsection{Input Distance Function}

\[
D_i(x,y)
=
\sup \{\lambda > 0 : x/\lambda \in V(y)\}.
\]

\medskip

The input distance function measures the maximal proportional contraction
of the input vector $x$ along the ray through $x$ while remaining feasible.

\medskip

Note that
\[
x \in V(y) \;\Longleftrightarrow\; D_i(x,y) \ge 1.
\]

\subsection*{Properties of the Input Distance Function}

\begin{itemize}
    \item[$D_i$.1] $D_i(x,0)=+\infty \text{ if } x \ne 0, D_i(0,y)=0.$
    
    \item[$D_i$.2] $D_i(x,y') \le D_i(x,y)$\\
    if $y \le y'$\\
    (non-increasing in $y$).

    \item[$D_i$.3] $D_i(x,y) \le D_i(x',y)$ \qquad if $x' \le x$\\
    (non-decreasing in $x$).

    \item[$D_i$.4] $D_i(\mu x,y)=\mu D_i(x,y), \qquad \mu>0$\\
    (homogeneous of degree 1 in $x$).

    \item[$D_i$.5] $D_i(\lambda x+(1-\lambda)x',y) \ge \lambda D_i(x,y)+(1-\lambda)D_i(x',y)$
    (concave in $x$).
\end{itemize}

\subsection*{Relationship Between Distance Functions}

In general, there is no direct relationship between the output distance
function and the input distance function.

\medskip

However, under constant returns to scale (CRS),
\[
D_o(x,y) = \frac{1}{D_i(x,y)}.
\]





\subsection*{Proof}

Under CRS, we have
\[
D_o(x,y)
=
\min \{\theta > 0 : y/\theta \in Y(x)\}.
\]

Using CRS,
\[
y/\theta \in Y(x)
\;\Longleftrightarrow\;
y \in \theta Y(x)
\;\Longleftrightarrow\;
y \in Y(\theta x).
\]

Thus,
\[
D_o(x,y)
=
\min \{\theta > 0 : \theta x \in V(y)\}.
\]

This implies
\[
D_o(x,y)
=
\frac{1}{\max \{ \lambda > 0 : x/\lambda \in V(y)\}}
=
\frac{1}{D_i(x,y)}.
\quad \blacksquare
\]

\subsection*{A Digression: Directional Distance Function}

Following Luenberger (1992), Chambers et al. (1996), and Chambers (2002),
the directional technology distance function is defined as
\[
D^{\rightarrow}_T(x,y;g_x,g_y)
=
\max \{\beta \in \mathbb{R} : (x - \beta g_x,\ y + \beta g_y) \in T\},
\]
where
\[
g_x \in \mathbb{R}^N,\quad g_y \in \mathbb{R}^M,\quad (g_x,g_y) \ne (0_N,0_M).
\]

If no such $\beta$ exists, then
\[
D^{\rightarrow}_T(x,y;g_x,g_y) = -\infty.
\]

\medskip

Here $(g_x,g_y)$ is a reference direction for inputs and outputs.
The function $D^{\rightarrow}_T(x,y;g_x,g_y)$ represents the maximal translation
of $(x,y)$ in the direction
\[
(-g_x,\ g_y)
\]
that keeps the point inside the technology set $T$.

\medskip

A firm is said to be \textbf{efficient} if
\[
D^{\rightarrow}_T(x,y;g_x,g_y) = 0.
\]

\medskip

The directional distance function allows both inputs and outputs to change
simultaneously.

\medskip

\textbf{Remark.}  
The output and input distance functions are special cases of the directional
distance function.

Moreover, there are important dual relationships between the directional
distance function and the profit function.




\section{Indirect Objective Functions}

\subsection{Cost Function}

The cost function is defined as
\[
c(w,y)
=
\min_{x \ge 0} \{ w \cdot x : x \in V(y) \},
\qquad w \in \mathbb{R}^N_+.
\]

The solution is denoted by
\[
x(w,y),
\]
which is called the \textbf{conditional factor demand function}.

\subsection*{Properties of the Cost Function}

\begin{itemize}
    \item[c.1] $c(w,y) > 0$ for $y>0$ \\
    (nonnegativity)

    \item[c.2] If $w' \ge w$, then $c(w',y) \ge c(w,y)$\\
    (nondecreasing in $w$)

    \item[c.3] $c(tw+(1-t)w',y)$\\
    $\ge t c(w,y)+(1-t)c(w',y),$\\
    $0<t<1$\\
    (concave in $w$)

    \item[c.4] $c(tw,y)=t c(w,y),$ $t>0$\\
    (homogeneous of degree 1 in $w$)

    \item[c.5] If $y \ge y'$, then $c(w,y) \ge c(w,y')$\\
    (nondecreasing in $y$)

    \item[c.6] $c(w,0)=0$\\
    (no fixed cost; not required in general)

    \item[c.7] If $c(w,y)$ is differentiable at $(w,y)$ and $w_i>0$,\\
    then $x_i(w,y)=\frac{\partial c(w,y)}{\partial w_i}, $\\
    $i=1,\dots,N.$\\
    (\textbf{Shephard’s Lemma})
\end{itemize}

\subsection*{Proof}

Property C.1 is immediate.

\medskip

For C.2, let $x$ and $x'$ be the cost-minimizing input bundles under prices $w$ and $w'$, respectively. Then
\[
w \cdot x \leq w \cdot x'.
\]
Since $w \leq w'$ (componentwise), we also have
\[
w \cdot x' \leq w' \cdot x'.
\]
Therefore,
\[
c(w,y) = w \cdot x \leq w' \cdot x' = c(w',y).
\]

\medskip

For C.3, let
\[
w'' = t w + (1-t) w', \quad 0 \leq t \leq 1.
\]
Let $x''$ be the cost-minimizing bundle at $w''$. Then
\[
c(w'',y) = w'' \cdot x'' = t\, w \cdot x'' + (1-t)\, w' \cdot x''.
\]
By definition of the cost function,
\[
w \cdot x'' \geq c(w,y), \quad w' \cdot x'' \geq c(w',y).
\]
Hence,
\[
c(w'',y) \geq t\, c(w,y) + (1-t)\, c(w',y).
\]

\medskip

For C.4, let $x$ be a cost-minimizing bundle at $w$. Then we claim that $x$ also minimizes cost at $tw$ for $t>0$. Suppose not, and let $x'$ be a minimizing bundle at $tw$ such that
\[
t\, w \cdot x' < t\, w \cdot x.
\]
Dividing both sides by $t>0$, we obtain
\[
w \cdot x' < w \cdot x,
\]
which contradicts the optimality of $x$ at $w$. Therefore,
\[
c(tw,y) = t\, c(w,y).
\]

\medskip

For C.5, if $y \leq y'$, then
\[
V(y) \subseteq V(y').
\]
Any feasible input bundle for producing $y$ is also feasible for producing $y'$. Therefore,
\[
c(w,y) \leq c(w,y').
\]

\medskip

For C.6, since producing zero output requires no inputs,
\[
c(w,0) = 0.
\]

\medskip

For C.7 (Shephard's Lemma), consider the Lagrangian
\[
\mathcal{L}(x,\lambda) = w \cdot x + \lambda \bigl( y - f(x) \bigr).
\]
By the envelope theorem,
\[
\frac{\partial c(w,y)}{\partial w_i} = x_i(w,y), \quad i=1,\dots,N.
\]
\hfill $\blacksquare$

\subsection*{Continuity}

$c(w,y)$ is continuous if $V(y)$ is strictly convex.

\subsection*{Comparative Statics}


The cost function is concave in $w$. Consider the Hessian matrix of $c(w,y)$ with respect to $w$:
\[
D_w^2 c(w,y)
=
\left[
\frac{\partial^2 c(w,y)}{\partial w_i \partial w_j}
\right]_{i,j=1}^N.
\]

By Shephard's lemma,
\[
x_i(w,y) = \frac{\partial c(w,y)}{\partial w_i}, \quad i=1,\dots,N.
\]

Hence, the Hessian can be written as
\[
D_w^2 c(w,y)
=
\left[
\frac{\partial x_i(w,y)}{\partial w_j}
\right]_{i,j=1}^N.
\]

Explicitly,
\[
D_w^2 c(w,y)
=
\begin{pmatrix}
\frac{\partial x_1}{\partial w_1} & \frac{\partial x_1}{\partial w_2} & \cdots & \frac{\partial x_1}{\partial w_N} \\
\frac{\partial x_2}{\partial w_1} & \frac{\partial x_2}{\partial w_2} & \cdots & \frac{\partial x_2}{\partial w_N} \\
\vdots & \vdots & \ddots & \vdots \\
\frac{\partial x_N}{\partial w_1} & \frac{\partial x_N}{\partial w_2} & \cdots & \frac{\partial x_N}{\partial w_N}
\end{pmatrix}.
\]

\textbf{A. Symmetry (Young’s theorem)}

\[
\frac{\partial x_i(w,y)}{\partial w_j}
=
\frac{\partial x_j(w,y)}{\partial w_i}
\]

Cross-price effects are symmetric.

\medskip

\textbf{B. Own-price effect}

\[
\frac{\partial x_i(w,y)}{\partial w_i} \le 0
\]

Own-price effects are nonpositive.



\textbf{C. Fundamental inequality of cost minimization}

\[
(x^1 - x^2)\cdot (w^1 - w^2) \le 0.
\]

Indeed, since $x^1$ minimizes cost at $w^1$ and $x^2$ minimizes cost at $w^2$,
we have
\[
w^1 \cdot x^1 \le w^1 \cdot x^2,
\qquad
w^2 \cdot x^2 \le w^2 \cdot x^1.
\]
Adding these two inequalities and rearranging yields
\[
(x^1 - x^2)\cdot (w^1 - w^2) \le 0.
\]

\begin{tcolorbox}[breakable]

\subsection*{Exercise 2-2}

Derive a cost function that is consistent with the CES production function
\[
y = \left(a x_1^{\rho} + b x_2^{\rho}\right)^{1/\rho},
\qquad a>0,\; b>0.
\]

\end{tcolorbox}

\begin{tcolorbox}[breakable]

\subsection*{Exercise 2-3}

Suppose the cost function is
\[
c(w,y)=w_2\left[1+y+\ln\left(\frac{w_1}{w_2}\right)\right],
\qquad y \in \mathbb{R}_+.
\]

Derive the corresponding production function.

\end{tcolorbox}

\begin{tcolorbox}[breakable]

\subsection*{Exercise 2-4}

If the technology has nonincreasing, constant, or nondecreasing returns to scale,
show that
\[
c(w,ty) \le t\,c(w,y)
\]
for $0<t<1$, for all $t$, or for all $t>1$, respectively.

If y is one-dimensional and there are constant returns to scale, show that
\[
c(w,y)=y\,c(w,1).
\]

\end{tcolorbox}

\begin{tcolorbox}[breakable]

\subsection*{Exercise 2-5}

Consider a single-output production function
\[
y=f(x), \qquad x \in \mathbb{R}_+^N.
\]

The elasticity of scale is defined as
\[
\mathcal{E}(x,y)
=
\sum_{i=1}^N \frac{\partial f(x)}{\partial x_i}\frac{x_i}{y}
=
\sum_{i=1}^N \epsilon_i.
\]
That is, the elasticity of scale is the sum of the individual scale elasticities.

Let the corresponding cost function be $c(w,y)$.

Define the elasticity of size by
\[
\frac{1}{\mathcal{E}^*(w,y)}
=
\frac{\partial c(w,y)}{\partial y}\frac{y}{c(w,y)}.
\]

Show that
\[
\mathcal{E}(x,y)=\mathcal{E}^*(w,y)
\]
when cost-minimizing inputs are chosen.

\end{tcolorbox}




\begin{tcolorbox}[breakable]

\subsection*{Exercise 2-6}

Suppose
\[
c(w,y)=y\,c(w)
\]
is the cost function, and $p$ is the output price.

If the output market is competitive, then producer profit is zero, so
\[
p=c(w).
\]

Now output is not fixed but determined by a demand schedule
\[
D(p)=y.
\]

Derive the response of input demand $x_i$ to a change in input price $w_j$,
or in elasticity form
\[
\frac{\partial x_i}{\partial w_j} \frac{w_j}{x_i},
\].

If the effect can be decomposed into multiple parts, show the decomposition.

\end{tcolorbox}



\subsection{Revenue Function}

The revenue function is defined as
\[
R(p,x)
=
\max_{y \ge 0} \{p \cdot y : y \in Y(x)\},
\qquad p \in \mathbb{R}_+^M.
\]

The solution is denoted by
\[
y(p,x).
\]

\subsection*{Properties of the Revenue Function}

\begin{itemize}
    \item[R.1] $R(p,x)>0$\\
    (nonnegativity)

    \item[R.2] If $p \le p'$, then $R(p,x) \le R(p',x)$\\
    (nondecreasing in $p$)

    \item[R.3] $R(tp+(1-t)p',x) \le tR(p,x)+(1-t)R(p',x),$\\
    $0\le t\le 1$\\
    (convex in $p$)

    \item[R.4] $R(tp,x)=tR(p,x),$\\
    $t>0$\\
    (homogeneous of degree 1 in $p$)

    \item[R.5] If $x' \ge x$, then $R(p,x') \ge R(p,x)$\\
    (nondecreasing in $x$)

    \item[R.6] If $R(p,x)$ is differentiable in $p$, then a unique
    revenue-maximizing output vector exists, and
    \[
    y_i(p,x)=\frac{\partial R(p,x)}{\partial p_i},
    \qquad i=1,\dots,M.
    \]
    This is the \textbf{Samuelson--McFadden lemma}.
\end{itemize}


\subsection*{Comparative Statics}

Since the revenue function is convex in $p$, its Hessian with respect to $p$,
\[
D^2_{pp}R(p,x)
=
\nabla_p y(p,x),
\]
is positive semidefinite and symmetric.

\medskip

\textbf{A. Symmetry}

\[
\frac{\partial y_i(p,x)}{\partial p_j}
=
\frac{\partial y_j(p,x)}{\partial p_i}
\]
by Young's theorem.

\medskip

\textbf{B. Own-price effects}

\[
\frac{\partial y_i(p,x)}{\partial p_i} \ge 0.
\]

Thus, own-price effects are nonnegative.

\medskip

\textbf{C. Fundamental inequality of revenue maximization}

\[
(y^1-y^2)\cdot (p^1-p^2) \ge 0.
\]

Indeed, since $y^1$ maximizes revenue at $p^1$ and $y^2$ maximizes revenue at $p^2$,
we have
\[
p^1 \cdot y^1 \ge p^1 \cdot y^2,
\qquad
p^2 \cdot y^2 \ge p^2 \cdot y^1.
\]
Adding the two inequalities and rearranging gives
\[
(y^1-y^2)\cdot (p^1-p^2) \ge 0.
\]

\subsection*{Example 2}

Consider the output set
\[
Y(x)=\{y : a_1 y_1 \le x,\ a_2 y_2 \le x,\ a_1,a_2>0\},
\]
which represents a Leontief technology.

This output set satisfies R.1--R.5.
For any $p>0$, revenue is maximized at the vertex where
\[
a_i y_i = x
\qquad \forall i.
\]

Hence the associated revenue function is
\[
R(p,x)
=
x \sum_{i=1}^M \frac{p_i}{a_i}.
\]

This revenue function is differentiable in $p$, and therefore
\[
y_i(p,x)
=
\frac{\partial R(p,x)}{\partial p_i}
=
\frac{x}{a_i}.
\]

\subsection*{Example 3}

Consider the output set
\[
Y(x)=\left\{y : \sum_{i=1}^M \beta_i y_i \le x,\ \beta_i>0\right\},
\]
which represents a linear technology.



Then the revenue function is
\[
R(p,x)
=
\max \left\{ \frac{p_1}{\beta_1}, \dots, \frac{p_M}{\beta_M} \right\} \cdot x.
\]

This revenue function is not differentiable.
If
\[
\frac{p_i}{\beta_i} = \frac{p_j}{\beta_j}
\]
for some $i \ne j$, then the maximizing output vector is not unique.

\section{Profit Function}

The profit function is defined as
\[
\Pi(p,w)
=
\max_{x,y} \{p\cdot y - w\cdot x : (x,y)\in T,\ p>0,\ w>0\}.
\tag{1}
\]

Equivalently,
\[
\Pi(p,w)
=
\max_{y} \{p\cdot y - c(w,y)\}.
\tag{2}
\]

Also,
\[
\Pi(p,w)
=
\max_{x} \{R(p,x) - w\cdot x\}.
\tag{3}
\]

\subsection*{Interpretation}

Expression (2) implies a two-stage maximization:

\begin{itemize}
    \item First, find the cost-minimizing input bundle for each $y$.
    \item Second, choose the profit-maximizing output level.
\end{itemize}

Similarly, (3) implies:

\begin{itemize}
    \item First, find the revenue-maximizing output for each $x$.
    \item Second, choose the profit-maximizing input level.
\end{itemize}

\subsection*{Proof of (1) $\Leftrightarrow$ (2)}

We need to show that
\[
x(p,w) = x(w, y(p,w)),
\]
i.e., the profit-maximizing input equals the cost-minimizing input
for the profit-maximizing output.

Suppose $x(p,w)=x' \ne x(w,y(p,w))$.

Then
\[
\Pi(p,w)
=
p\cdot y(p,w) - w\cdot x'
\ge
p\cdot y(p,w) - w\cdot x(w,y(p,w)).
\]

This implies
\[
w\cdot x' \le w\cdot x(w,y(p,w)),
\]
which contradicts the definition of $x(w,y(p,w))$ as the
cost-minimizing input bundle unless
\[
x' = x(w,y(p,w)).
\]
\hfill $\blacksquare$

\begin{tcolorbox}

\subsection*{Exercise 7}

Prove $(1) \Leftrightarrow (3)$ by showing that
\[
y(p, x(p,w)) = y(p,w).
\]

\end{tcolorbox}

\begin{tcolorbox}

\subsection*{Exercise 8}

Prove that a profit-maximizing producer chooses output levels $y(p,w)$ such that
\[
p_i = \frac{\partial c(w,y)}{\partial y_i}.
\]

Moreover, show that the marginal cost
\[
\frac{\partial c(w,y)}{\partial y_i}
\]
increases in $y_i$ at the optimal choice.

\end{tcolorbox}

\subsection*{Properties of the Profit Function}

\begin{itemize}
    \item[$\Pi$.1] $\Pi(p,w)\ge 0$ and exists.
    
    \item[$\Pi$.2] If $p \ge p'$, then
    \[
    \Pi(p,w)\ge \Pi(p',w)
    \]
    \hfill (nondecreasing in $p$)
    
    \item[$\Pi$.3] If $w \ge w'$, then
    \[
    \Pi(p,w)\le \Pi(p,w')
    \]
    \hfill (nonincreasing in $w$)
    
    \item[$\Pi$.4] $\Pi(p,w)$ is convex in $(p,w)$.
    
    \item[$\Pi$.5]
    \[
    \Pi(tp,tw)=t\Pi(p,w), \qquad t>0
    \]
    \hfill (homogeneous of degree 1 in $(p,w)$)
    
    \item[$\Pi$.6] If the profit function is differentiable in $(p,w)$, then
    \[
    \frac{\partial \Pi(p,w)}{\partial p_i}=y_i(p,w),
    \qquad
    \frac{\partial \Pi(p,w)}{\partial w_i}=-x_i(p,w)
    \]
    \hfill (Hotelling's lemma)
\end{itemize}

\subsection*{Proof Sketch}

$\Pi(p,w)$ exists since $T$ is compact and $p\cdot y - w\cdot x$ is continuous.
The firm can always choose $(0_N,0_M)$ and earn zero profit, so
\[
\Pi(p,w)\ge 0.
\]

\medskip

If $(x',y')$ is the profit-maximizing choice associated with $(p',w)$, then
\[
\Pi(p',w)=p'\cdot y' - w\cdot x'.
\]
If $p\ge p'$, then
\[
\Pi(p',w)=p'\cdot y' - w\cdot x'
\le p\cdot y' - w\cdot x'
\le \Pi(p,w).
\]
Hence $\Pi$.2 holds.

\medskip

Similarly, if $(x,y)$ is the profit-maximizing choice associated with $(p,w)$, then
\[
\Pi(p,w)=p\cdot y - w\cdot x
\le p\cdot y - w'\cdot x
\le \Pi(p,w'),
\]
which proves $\Pi$.3.

\medskip

To show convexity, let $(x^2,y^2)$ be the profit-maximizing choice for
\[
(p^2,w^2)=\bigl(\theta p^0+(1-\theta)p^1,\ \theta w^0+(1-\theta)w^1\bigr),
\qquad 0\le \theta \le 1.
\]
Then
\[
\Pi(p^0,w^0)\ge p^0\cdot y^2 - w^0\cdot x^2,
\]
\[
\Pi(p^1,w^1)\ge p^1\cdot y^2 - w^1\cdot x^2.
\]
Therefore,
\[
\theta \Pi(p^0,w^0)+(1-\theta)\Pi(p^1,w^1)
\ge
\Pi(p^2,w^2),
\]
so $\Pi(p,w)$ is convex.

\medskip

For homogeneity,
\[
\Pi(tp,tw)
=
\max_y \{tp\cdot y - c(tw,y)\}
=
\max_y \{tp\cdot y - t c(w,y)\}
=
t\max_y \{p\cdot y - c(w,y)\}
=
t\Pi(p,w).
\]

\medskip

Applying the envelope theorem to
\[
\mathcal{L}=p\cdot y - w\cdot x + \lambda Y(x,y),
\]
we obtain
\[
\frac{\partial \Pi(p,w)}{\partial p_i}=y_i(p,w),
\qquad
\frac{\partial \Pi(p,w)}{\partial w_i}=-x_i(p,w).
\]
\hfill $\blacksquare$

\begin{tcolorbox}

\subsection*{Exercise 9}

Let the production function be CES:
\[
y = \left(x_1^\rho + x_2^\rho\right)^{\beta/\rho},
\qquad \beta<1,\quad 0\ne \rho<1.
\]

Derive the input demand functions and the profit function.


\end{tcolorbox}



\subsection*{Comparative Statics}

By homogeneity and Hotelling's lemma,
\[
y_i(tp,tw)=y_i(p,w),
\qquad
x_i(tp,tw)=x_i(p,w).
\]

Since $\Pi(p,w)$ is convex in $q=(p,w)$, the Hessian
\[
\nabla_{qq}^2 \Pi(p,w)
\]
is positive semidefinite.

\medskip

In particular,
\[
\frac{\partial y_i(p,w)}{\partial p_i} > 0,
\qquad
\frac{\partial x_i(p,w)}{\partial w_i} < 0
\]
\hfill (own-price effects)

and symmetry implies
\[
\frac{\partial x_i(p,w)}{\partial w_j}
=
\frac{\partial x_j(p,w)}{\partial w_i},
\]
\[
\frac{\partial y_i(p,w)}{\partial p_j}
=
\frac{\partial y_j(p,w)}{\partial p_i},
\]
\[
\frac{\partial y_i(p,w)}{\partial w_j}
=
-\frac{\partial x_j(p,w)}{\partial p_i}.
\]


\subsection*{(continued)}

\textbf{Fundamental inequality of profit maximization}

\[
(p^1-p)\cdot (y^1-y) - (w^1-w)\cdot (x^1-x) \ge 0.
\]

Indeed, profit maximization implies
\[
p\cdot y - w\cdot x \ge p\cdot y^1 - w\cdot x^1,
\]
and
\[
p^1 \cdot y^1 - w^1 \cdot x^1 \ge p^1 \cdot y - w^1 \cdot x.
\]
Adding these inequalities and rearranging yields the result above.

\subsection*{Decomposition of Derived Demands and Output Supplies}

We know that
\[
y(p,x(p,w)) = y(p,w),
\qquad
x(w,y(p,w)) = x(p,w).
\]

Therefore,
\[
\frac{\partial y_i(p,w)}{\partial p_j}
=
\frac{\partial y_i(p,x)}{\partial p_j}
+
\sum_{v=1}^N
\frac{\partial y_i(p,x)}{\partial x_v}
\frac{\partial x_v(p,w)}{\partial p_j}.
\]

This may be interpreted as

\[
\text{movement along the boundary of } Y(x)
+
\text{a shift in } Y(x).
\]

Similarly,
\[
\frac{\partial y_i(p,w)}{\partial w_k}
=
\sum_{v=1}^N
\frac{\partial y_i(p,x)}{\partial x_v}
\frac{\partial x_v(p,w)}{\partial w_k}.
\]

Also,
\[
\frac{\partial x_i(p,w)}{\partial p_j}
=
\sum_{v=1}^M
\frac{\partial x_i(w,y)}{\partial y_v}
\frac{\partial y_v(p,w)}{\partial p_j},
\]

and
\[
\frac{\partial x_i(p,w)}{\partial w_k}
=
\frac{\partial x_i(w,y)}{\partial w_k}
+
\sum_{v=1}^M
\frac{\partial x_i(w,y)}{\partial y_v}
\frac{\partial y_v(p,w)}{\partial w_k}.
\]

This may be interpreted as

\[
\text{movement along the boundary of } V(y)
+
\text{a shift in } V(y).
\]

\subsection*{Example 4}

Consider
\[
Y(x)=\left\{y : a_1 y_1 + a_2 y_2 \le x^{1/2},\ a_1,a_2>0\right\}.
\]

For given $x$, consider a revenue-maximization problem.
Assume that
\[
p_1 > \frac{a_1}{a_2} p_2.
\]

Then the optimal solution is a corner solution:
\[
y_2 = 0,
\qquad
y_1 = \frac{x^{1/2}}{a_1}.
\]

The associated profit-maximization problem becomes
\[
\max_x \left( \frac{p_1}{a_1} x^{1/2} - wx \right).
\]

The first-order condition implies
\[
x^{-1/2} = 2a_1 \left(\frac{w}{p_1}\right).
\]

Now suppose $p_1$ rises. Again, $y_2$ is not produced, but the
profit-maximizing input $x$ increases, so the optimal production point
moves to a new point on the boundary.

In this case,
\[
\frac{\partial y_i(p,w)}{\partial p_j}
=
\frac{\partial y_i(p,x)}{\partial p_j}
+
\sum_{v=1}^N
\frac{\partial y_i(p,x)}{\partial x_v}
\frac{\partial x_v(p,w)}{\partial p_j}.
\]

\subsection*{Le Chatelier--Samuelson Principle}

Let
\[
y^* = y(p,w)
\]
be the point where the short-run optimum
(no shift in $V(y)$) and the long-run optimum coincide.

We know that
\[
\frac{\partial x_i(p,w)}{\partial w_k}
=
\frac{\partial x_i(w,y^*)}{\partial w_k}
+
\sum_{v=1}^M
\frac{\partial x_i(w,y)}{\partial y_v}
\frac{\partial y_v(p,w)}{\partial w_k}.
\]

By symmetry,
\[
\frac{\partial y_v(p,w)}{\partial w_k}
=
-
\frac{\partial x_k(p,w)}{\partial p_v}.
\]

Combining these gives
\[
\frac{\partial x_i(p,w)}{\partial w_k}
=
\frac{\partial x_i(w,y^*)}{\partial w_k}
-
\sum_{v=1}^M \sum_{s=1}^M
\frac{\partial x_i(w,y)}{\partial y_v}
\frac{\partial y_s(p,w)}{\partial p_v}
\frac{\partial x_k(w,y)}{\partial y_s}.
\]

In particular, when $i=k$, this expression compares the long-run and short-run
responses of input demand to a change in input price.

This is the essence of the \textbf{Le Chatelier--Samuelson principle}:
the long-run response is at least as large in magnitude as the short-run response.




\subsection*{(continued)}

Notice that
\[
\sum_{v=1}^M \sum_{s=1}^M
\frac{\partial x_i(w,y)}{\partial y_v}
\frac{\partial y_s(p,w)}{\partial p_v}
\frac{\partial x_i(w,y)}{\partial y_s}
\]
is a quadratic form, which can be written as
\[
\left( \nabla_y x_i(w,y) \right)^\top
\left( \nabla^2_{pp} \Pi(p,w) \right)
\left( \nabla_y x_i(w,y) \right).
\]

Since the profit function $\Pi(p,w)$ is convex in $(p,w)$,
its Hessian $\nabla^2_{pp}\Pi(p,w)$ is positive semidefinite.
Therefore, the quadratic form above is nonnegative.

\medskip

Thus,
\[
\frac{\partial x_i(p,w)}{\partial w_i}
\le
\frac{\partial x_i(w,y^*)}{\partial w_i}.
\]

Similarly, we can show
\[
\frac{\partial y_i(p,w)}{\partial p_i}
\ge
\frac{\partial y_i(p,x^*)}{\partial p_i}.
\]

\medskip

\textbf{Conclusion (Le Chatelier--Samuelson Principle):}

The long-run response to a price change is more elastic than the short-run response.

\[
|\text{long-run effect}| \ge |\text{short-run effect}|.
\]

\section{3 Principle of Duality}

Indirect objective functions such as the cost, revenue, and profit functions
are defined using the primal technology.

For example, the input requirement set $V(y)$ is used to define the cost function:
\[
c(w,y)=\min_{x \ge 0} \{ w \cdot x : x \in V(y) \}.
\]

\medskip

If we impose restrictions on the primal technology, we implicitly restrict
the form of the dual functions.

\medskip

The natural question is:

\begin{quote}
Can we reverse the process?  
Can we use observed data on dual functions to recover the underlying technology?
\end{quote}

\medskip

\textbf{Principle of Duality:}

Under certain conditions, this is indeed possible.

Specifying an indirect objective function (such as $c(w,y)$, $R(p,x)$, or $\Pi(p,w)$)
is equivalent to specifying the underlying primal technology.




\subsection*{3.1 Illustration of the Duality}

Generate a half-space at a particular $w$ and $y$:
\[
N(w,y)=\{x : w\cdot x \le c(w,y)\}.
\]

That is, the cost function defines a half-space.

\medskip

In Figure 12, the half-space $N(w,y)$ is the area below the line $AB$.
Any point below the hyperplane
\[
w\cdot x = c(w,y)
\]
cannot produce $y$.
Therefore, the half-space above $AB$ should contain $V(y)$.

Now consider the half-space defined by
\[
N(w,y)=\{x : w\cdot x \ge c(w,y)\}.
\]

If the hyperplane
\[
w\cdot x = c(w,y)
\]
is represented by $CD$ in Figure 12, then the area above $CD$ must contain $V(y)$.
Hence $V(y)$ must lie above both $CD$ and $AB$.

If this construction is repeated for all possible price vectors,
the lower boundary of the intersection of these half-spaces takes a shape
similar to that of convex isoquants.

\medskip

The constructed input requirement set is defined as
\[
V^*(y)=\{x : w\cdot x \ge c(w,y)\ \forall w>0\}.
\]

The question is:
\[
V^*(y)=V(y)\, ?
\]

\subsection*{Properties of $V^*(y)$}

\begin{itemize}
    \item[$V^*$.1] If $x \in V^*(y)$, then for any $x' \ge x$,
    \[
    x' \in V^*(y)
    \]
    \hfill (monotonicity)

    \item[$V^*$.2]
    \[
    0_N \in V^*(0_M),
    \qquad
    0_N \notin V^*(y)\ \text{for } y \ne 0
    \]

    \item[$V^*$.3] $V^*(y)$ is always convex, regardless of whether or not $V(y)$ is convex.

    \item[$V^*$.4] $V^*(y)$ is closed and nonempty.
\end{itemize}




\subsection*{(continued)}

\subsection*{Proof of Properties of $V^*(y)$}

Since $x \in V^*(y)$, we have
\[
w \cdot x \ge c(w,y).
\]
If $x' \ge x$, then
\[
w \cdot x' \ge w \cdot x \ge c(w,y),
\]
so $x' \in V^*(y)$. This proves $V^*.1$.

\medskip

Since $c(w,0)=0$, we have
\[
w \cdot 0_N = c(w,0),
\]
so $0_N \in V^*(0_M)$.

On the other hand, for any $y \ne 0$, $c(w,y)>0$, so
\[
w \cdot 0_N = 0 < c(w,y),
\]
and hence $0_N \notin V^*(y)$. This proves $V^*.2$.

\medskip

Suppose $x^0, x' \in V^*(y)$, so
\[
w \cdot x^0 \ge c(w,y),
\qquad
w \cdot x' \ge c(w,y).
\]
Let
\[
x^2 = \theta x^0 + (1-\theta)x', \qquad 0 \le \theta \le 1.
\]
Then
\[
w \cdot x^2
=
\theta w \cdot x^0 + (1-\theta) w \cdot x'
\ge
\theta c(w,y) + (1-\theta)c(w,y)
=
c(w,y).
\]
Thus $x^2 \in V^*(y)$, so $V^*(y)$ is convex. This proves $V^*.3$.

\medskip

$V^*.4$ is immediate. \hfill $\blacksquare$

\subsection*{Interpretation}

Property $V^*.3$ implies that $V^*(y)$ is the \textbf{convex hull} of $V(y)$.

That is,
\[
V^*(y) \supseteq V(y),
\qquad
V^*(y)=V(y) \quad \text{if $V(y)$ is convex}.
\]

\medskip

\subsection*{Economic Interpretation}

Consider a nonconvex $V(y)$ as in Figure 13.

If the price line is $AB$ or $DE$, then an input bundle on those segments is chosen.

If the price line is $BC$ (or $CD$), then a corner point such as $D$ (or $B$)
is chosen instead.

Thus, points like $C$ in the nonconvex region of $V(y)$ are never chosen
in cost minimization.

\medskip

Therefore, although
\[
V^*(y) \ne V(y),
\]
no economically relevant information is lost by replacing $V(y)$ with $V^*(y)$.

\begin{tcolorbox}

\subsection*{Exercise 10}

Given a profit function $\Pi(p,w)$, define a set...


Given a profit function $\Pi(p,w)$, define the set
\[
T^*
=
\{(x,y)\in \mathbb{R}_+^{N+M} : p\cdot y - w\cdot x \le \Pi(p,w)
\ \text{for all } p>0,\ w>0\}.
\]

Derive the properties of $T^*$ and compare them with those of $T$.

\end{tcolorbox}

\section{4 Primal vs.\ Dual: Some Dualities}

We can establish the following dual relationships between the input
distance function and the cost function:

\[
(1)\qquad
c(w,y)
=
\min \{w\cdot x : D_i(x,y)\ge 1\},
\]

\[
(2)\qquad
c(w,y)
=
\min_x \left\{ \frac{w\cdot x}{D_i(x,y)} \right\},
\]

\[
(3)\qquad
D_i(x,y)
=
\min \{ w\cdot x : c(w,y)\ge 1\},
\]

\[
(4)\qquad
D_i(x,y)
=
\min_w \left\{ \frac{w\cdot x}{c(w,y)} \right\}.
\]

\subsection*{Proof Sketch}

Relation (1) is immediate from the definition, since
\[
D_i(x,y)\ge 1
\quad \Longleftrightarrow \quad
x \in V(y).
\]

Now,
\[
D_i(x,y)
=
\max \left\{ \lambda : \frac{x}{\lambda}\in V(y) \right\}.
\]

Since
\[
\frac{x}{\lambda}\in V(y)
\quad \Longleftrightarrow \quad
w\cdot \frac{x}{\lambda} \ge c(w,y)
\quad \forall w>0,
\]
we obtain
\[
D_i(x,y)
=
\max \left\{ \lambda : \lambda \le \frac{w\cdot x}{c(w,y)}
\quad \forall w>0 \right\}.
\]

Therefore,
\[
D_i(x,y)
=
\min_w \left\{ \frac{w\cdot x}{c(w,y)} \right\},
\]
which proves (4).

\medskip

To prove (2) and (3), we use the following theorem.

\subsection*{Theorem 1}

Suppose $f(y,x)$ and $g(z,x)$ are real-valued functions, and both are
linearly homogeneous in $x\in X$.

Then
\[
\min_x \left\{
\frac{f(y,x)}{g(z,x)} : y\in Y,\ x\in X,\ z\in Z
\right\}
=
\min_x \left\{
f(y,x) : g(z,x)\ge 1,\ y\in Y,\ x\in X,\ z\in Z
\right\},
\]
whenever the minimum on each side exists.

\subsection*{Proof}

Let
\[
A
=
\left\{
\frac{f(y,x)}{g(z,x)} : y\in Y,\ x\in X,\ z\in Z
\right\},
\]
\[
B
=
\left\{
f(y,x) : g(z,x)\ge 1,\ y\in Y,\ x\in X,\ z\in Z
\right\},
\]
and define
\[
\alpha = \min A,
\qquad
\beta = \min B.
\]

We show that $\alpha=\beta$.

\medskip

First, let $b\in B$. Then for some $(y_0,x_0,z_0)$,
\[
b=f(y_0,x_0),
\qquad
g(z_0,x_0)\ge 1.
\]
Hence
\[
\frac{f(y_0,x_0)}{g(z_0,x_0)} \le f(y_0,x_0)=b,
\]
so there exists $a\in A$ such that $a\le b$.
Therefore,
\[
\alpha \le \beta.
\]

\medskip

Conversely, let $a\in A$. Then for some $(y_1,x_1,z_1)$,
\[
a = \frac{f(y_1,x_1)}{g(z_1,x_1)}.
\]
Define
\[
\tilde x = \frac{x_1}{g(z_1,x_1)}.
\]
By linear homogeneity,
\[
f(y_1,\tilde x)
=
\frac{f(y_1,x_1)}{g(z_1,x_1)},
\qquad
g(z_1,\tilde x)=1.
\]
Thus
\[
a=f(y_1,\tilde x)\in B.
\]
So
\[
A \subseteq B
\quad \Rightarrow \quad
\alpha \ge \beta.
\]

Hence
\[
\alpha=\beta.
\qquad \blacksquare
\]

\medskip

By this theorem, (1) implies (2), and (4) implies (3).


\begin{tcolorbox}

\subsection*{Exercise 11}

Consider the input distance function
\[
D_i(x,y)=\frac{x}{\sqrt{y_1^2+y_2^2}}.
\]

Derive the conditional factor demand functions and the cost function
corresponding to this input distance function.


\end{tcolorbox}

\subsection*{Dual Relations for Output Distance Function}

By symmetry with the input-distance case, we obtain:

\[
(5)\qquad
R(p,x)
=
\max \{ p\cdot y : D_o(x,y)\le 1 \},
\]

\[
(6)\qquad
R(p,x)
=
\max_y \left\{ \frac{p\cdot y}{D_o(x,y)} \right\},
\]

\[
(7)\qquad
D_o(x,y)
=
\max \{ p\cdot y : R(p,x)\le 1 \},
\]

\[
(8)\qquad
D_o(x,y)
=
\max_p \left\{ \frac{p\cdot y}{R(p,x)} \right\}.
\]

\subsection*{Geometric Interpretation}

For any output bundle $y$ that can be produced with input $x$,
there exists a normalized price vector $p$ such that the price line
is tangent to the production possibility frontier (PPF).

Let $A$ be such a point and $p^A$ the corresponding price vector:
\[
R(p^A,x)=p^A \cdot y^A = 1.
\]

Now consider another bundle $B$.

Let $C$ be the intersection between the ray $OB$ and the price line
associated with $p^A$.

Define
\[
\theta
=
\frac{OB}{OC}
=
\frac{p^A \cdot y^B}{p^A \cdot y^A}
=
p^A \cdot y^B.
\]

We want to find the price vector $p$ that maximizes $\theta$.

The optimal solution is given by some price vector $p^D$.

\medskip

\textbf{Result:}

The maximized value of $\theta$ is exactly equal to the output distance function:
\[
D_o(x,y)=\max_p \frac{p\cdot y}{R(p,x)}.
\]



\subsection*{(continued)}

Thus,
\[
\max \theta = OD = \frac{OB}{OC} = p^D \cdot y^B = D_o(x,y).
\]

This explains why
\[
D_o(x,y)=\max_p \{p\cdot y : R(p,x)=1\}.
\]

\begin{tcolorbox}

\subsection*{Exercise 12}

Illustrate (7) using a diagram similar to Figure 14.

\end{tcolorbox}

\subsection*{Example 5: Shadow Prices of Nonmarket Goods}

Using (5), construct the Lagrangian:
\[
\mathcal{L}=p\cdot y + \lambda(1 - D_o(x,y)).
\]

The solutions are $y(p,x)$ and $\lambda(p,x)$.

Applying the envelope theorem:
\[
\nabla_x R(p,x)
=
-\lambda(p,x)\nabla_x D_o(x,y(p,x)),
\]
\[
\nabla_p R(p,x)=y(p,x).
\]

The first-order condition is
\[
p = \lambda \nabla_y D_o(x,y).
\]

Thus,
\[
p\cdot y(p,x)
=
\lambda(p,x)\bigl[\nabla_y D_o(x,y)\cdot y\bigr]
=
\lambda(p,x) D_o(x,y)
\]
(by homogeneity of degree 1).

If we restrict attention to efficient observations where
\[
D_o(x,y)=1,
\]
then
\[
p\cdot y(p,x)=R(p,x)=\lambda(p,x).
\]

Substituting back:
\[
p = R(p,x)\nabla_y D_o(x,y).
\]

\medskip

\textbf{Application: Shadow Prices}

Suppose there are two outputs (e.g., electricity $y_1$ and heat $y_2$).

If $p_1$ is observed but $p_2$ is not, then
\[
\frac{p_2}{p_1}
=
\frac{\partial D_o(x,y)/\partial y_2}
{\partial D_o(x,y)/\partial y_1},
\]

which is the marginal rate of transformation (MRT).

Thus, estimating a differentiable output distance function allows us
to recover shadow prices.

\section{5 Cost, Revenue and Profit}

We can derive a cost function from a profit function.

Define the implicit cost function:
\[
c^*(w,y)=\max_p \{p\cdot y - \Pi(p,w)\}.
\]

\subsection*{Properties}

For any $y_0$,
\[
\Pi(p,w)\ge p\cdot y_0 - c(w,y_0),
\]
so
\[
c(w,y_0)\ge p\cdot y_0 - \Pi(p,w).
\]

Thus, $c(w,y_0)$ provides an upper bound.

Equality holds when $p$ is such that $y_0$ is the profit-maximizing output.

\medskip

Using first-order conditions:
\[
\frac{\partial}{\partial p}\bigl(p\cdot y - \Pi(p,w)\bigr)
=
y - y(p,w).
\]

Thus, the maximum occurs at
\[
y=y(p,w),
\]
and
\[
c^*(w,y)=p\cdot y - \Pi(p,w).
\]

\subsection*{Homogeneity}

\[
c^*(tw,y)
=
\max_p \{p\cdot y - \Pi(p,tw)\}
=
\max_p \{p\cdot y - t\Pi(p/t,w)\}
=
t\max_p \{(p/t)\cdot y - \Pi(p/t,w)\}
=
t c^*(w,y).
\]

\subsection*{Convexity}

Since $-\Pi(p,w)$ is concave in $w$, $c^*(w,y)$ is concave in $w$.

Moreover, $c^*(w,y)$ is convex in $y$.

Let
\[
y^2=\theta y^0 + (1-\theta)y^1.
\]

Then
\[
c^*(w,y^0)=p^0\cdot y^0 - \Pi(p^0,w),
\]
\[
c^*(w,y^1)=p^1\cdot y^1 - \Pi(p^1,w),
\]
\[
c^*(w,y^2)=p^2\cdot y^2 - \Pi(p^2,w),
\]

which implies convexity in $y$.




\subsection*{(continued)}

Also,
\[
c^*(w,y^0) \ge p^2 \cdot y^0 - \Pi(p^2,w),
\qquad
c^*(w,y^1) \ge p^2 \cdot y^1 - \Pi(p^2,w).
\]

Therefore,
\[
c^*(w,y^2)
=
p^2 \cdot y^2 - \Pi(p^2,w)
\]
\[
=
p^2 \cdot \bigl(\theta y^0 + (1-\theta)y^1\bigr) - \Pi(p^2,w)
\]
\[
=
\theta \bigl(p^2 \cdot y^0 - \Pi(p^2,w)\bigr)
+
(1-\theta)\bigl(p^2 \cdot y^1 - \Pi(p^2,w)\bigr)
\]
\[
\le
\theta c^*(w,y^0) + (1-\theta)c^*(w,y^1).
\]

Hence,
\[
c^*(w,y^2)
\le
\theta c^*(w,y^0) + (1-\theta)c^*(w,y^1),
\]
which proves that $c^*(w,y)$ is convex in $y$.

\medskip

The convexity of $c^*(w,y)$ means that profit is maximized only in the region
where marginal cost is increasing.


\begin{tcolorbox}

\subsection*{Exercise 13}

Derive a revenue function from a known profit function $\Pi(p,w)$.

\end{tcolorbox}



